% !TeX root = ../main.tex
\section{Mathematische Grundlagen}\label{sec:grundlagen}
\subsection{Gruppen und zyklische Gruppen}
\todo{Gruppe mit Verknüpfung definieren: Abgeschlossenheit, Assoziativität,
 neutrales Element und inverse Elemente. Ein kleines Beispiel geben.
 Danach endliche zyklische Gruppe, Generator und Gruppenordnung erklären.}
% Beispiel für eine Definitionsumgebung:
% \begin{definition}[Zyklische Gruppe]
% Eine Gruppe ...
% \end{definition}
\subsection{Modulare Arithmetik und Untergruppen}
\todo{Rechnen modulo einer Zahl erklären. Für das spätere Beispiel eine
 Untergruppe von $\Z_p^*$ mit primer Ordnung $q$ und $q\mid(p-1)$ einführen.
 Gruppenelemente werden modulo $p$, Exponenten modulo $q$ berechnet.}
\subsection{Das diskrete Logarithmusproblem}
\todo{Für $\G=\langle g\rangle$ und $y=g^x$ das Problem der Bestimmung
 von $x$ definieren. Leichte Vorwärtsberechnung und angenommene Schwierigkeit
 der Umkehrung unterscheiden. Die Schwierigkeit hängt von Gruppe und Parametern ab.}
\subsection{Verwendete Notation}
\begin{center}
\begin{tabular}{ll}
\toprule
Symbol & Bedeutung \\
\midrule
$\G$, $q$, $g$ & Zyklische Gruppe, prime Ordnung, Generator \\
$x\in\Z_q$, $y=g^x$ & Geheimer und öffentlicher Schlüssel \\
$r\in\Z_q$, $t=g^r$ & Frischer Zufallswert und Commitment \\
$c\in\Z_q$, $s\in\Z_q$ & Challenge und Response \\
\bottomrule
\end{tabular}
\end{center}
\todo{Notation im gesamten Paper konsistent halten; abweichende Vorzeichen
 in Quellen ausdrücklich übersetzen. Alle später benötigten Begriffe hier erklären.}
